%% ma: L12-B08-0138
%% lop: 12
%% chuong: 2
%% bai: 8
%% dang: 12.8.23
%% loai: TN4
%% muc: VD
%% dap_an: "A"
%% nguon: "CĐ2.3.pdf giáo viên gửi 27/10/2026 (M07, Lê Minh Khôi – Nguyễn Anh Khôi, Chương II Chủ đề 3 Biểu thức tọa độ của các phép toán vectơ), Câu 53"
%% kien_thuc: "Bài 8 (góc giữa hai vectơ), lớp 10 (phương trình chứa căn)"
%% trang_thai: chinh-thuc
\begin{ex}
Trong không gian $Oxyz$, cho tam giác $ABC$ có tọa độ các đỉnh $A(1;1;1)$, $B(5;1;-2)$ và $C(a;5;1)$. Tìm giá trị $a>0$ biết rằng $\cos\widehat{BAC}=\dfrac{12}{25}$.
\choice{\True $a=4$}{$a=3$}{$a=5$}{$a=1$}
\loigiai{$\overrightarrow{AB}=(4;0;-3)$, $\overrightarrow{AC}=(a-1;4;0)$: $\dfrac{4(a-1)}{5\sqrt{(a-1)^2+16}}=\dfrac{12}{25}$.

$\Leftrightarrow5(a-1)=3\sqrt{(a-1)^2+16}\Leftrightarrow a-1=3\Leftrightarrow a=4$.}
\end{ex}
